% !PRINTED 157
\weberpage{157}
Diese Bemerkung giebt die Vorschrift, nach der die transponirten Substitutionen zusammengesetzt werden, die sich in dem Satze ausspricht:

\par
8.\enspace\gesperrt{Sind $A_1$, $B_1$ die Transponirten zu $A$, $B$, so ist $A_1 B_1$ die Transponirte zu $BA$. In Zeichen:}

\weq{24}{(BA)_1 = A_1 B_1.}

Wenn wir die Substitutionsformeln (1) mit einem unbestimmten Factor $\eta_k$ multipliciren und dann die Summe in Bezug auf $k$ nehmen, so folgt:

\weq{25}{\sum y_k \eta_k = \sum^i \sum^k x_i\, a_i^{(k)}\, \eta_k,}

oder wenn wir

\weq{26}{\sum^k a_i^{(k)} \eta_k = \xi_i}

setzen,

\weq{27}{\sum y_k \eta_k = \sum x_i \xi_i.}

Wenn wir also die Substitutionen (1) durch (4) darstellen, so sind die Formeln (26) der Ausdruck für die Substitution

\[
(\xi) = A_1(\eta),
\]

und wir können also noch den Satz aussprechen:

\par
9.\enspace\gesperrt{Sind $A$, $A_1$ transponirte Substitutionen von einander von der Dimension $n$, und sind $x$, $y$, $\xi$, $\eta$ vier Systeme von Variablen, die mit einander durch die Substitutionen}

\weq{28}{(y) = A(x), \quad (\xi) = A_1(\eta)}

\gesperrt{zusammenhängen, so besteht die Identität}

\weq{29}{y_1\eta_1 + y_2\eta_2 + \cdots y_n\eta_n = x_1\xi_1 + x_2\xi_2 + \cdots x_n\xi_n.}

Wenn zwei Reihen von Variablen mit

\[
\begin{gathered} x_1, x_2, \ldots x_n \\ \xi_1, \xi_2, \ldots \xi_n \end{gathered}
\]

gleichzeitig durch die Substitutionen (28) in zwei neue Reihen

\[
\begin{gathered} y_1, y_2, \ldots y_n \\ \eta_1, \eta_2, \ldots \eta_n \end{gathered}
\]

transformirt werden, so heissen die beiden Variablenreihen mit einander \gesperrt{contragradient}, und die Formeln (28) stellen zwei mit einander \gesperrt{contragradiente Transformationen} dar.
